%% kapa: exam=1 round=25 year=2014 period=1 form=A subject=civil_law subject_ko=민법 seq=13 number=13
%% source: Q-Net 감정평가사 기출문제 (cst00309) artlSeq=5101207 file=extracted/1008/5101207/2014년 제25회 감정평가사 1차시험 기출문제/1교시(A형).hwp
%% conversion: kordoc 4.24.51 -> tools/gen1.py (auto); review=proofread-sonnet; flags=none
\begin{question}{13}
첨부(添附)에 관한 설명으로 옳은 것은? (다툼이 있으면 판례에 의함)
\begin{choices}
  \item 부동산에 부합된 물건이 타인의 권원에 의하여 부합되었더라도 사실상 분리가 불가능하여 부동산의 구성부분이 되었다면, 그 물건의 소유권은 부동산의 소유자에게 귀속된다.
  \item 동산이 부동산에 부합되면 그 동산의 소유권자는 부합물 위에 지분권을 취득한다.
  \item 동산이 부동산에 부합되더라도 그 동산을 목적으로 하는 다른 권리는 소멸하지 않는다.
  \item 가공물의 소유권은 원칙적으로 원재료의 가공자에게 속한다.
  \item 동산간의 부합에 있어서 합성물의 소유권은 원칙적으로 주종의 구별 없이 가액에 비례하여 공유로 한다.
\end{choices}
\end{question}
\answer{1}
